%% GENERATED by build_paper.py from the chapter .src.md files. Do not edit by hand.
\documentclass[12pt,a4paper]{report}
\usepackage[T1]{fontenc}
\usepackage[utf8]{inputenc}
\usepackage{mathptmx}
\usepackage[margin=2.5cm]{geometry}
\usepackage{setspace}
\onehalfspacing
\usepackage{csquotes}
\usepackage{microtype}
\usepackage[style=ieee,backend=biber]{biblatex}
\addbibresource{references.bib}
\usepackage{titlesec}
\titleformat{\chapter}[display]{\normalfont\large\bfseries\centering}{CHAPTER \thechapter}{0.3em}{\MakeUppercase}
\titlespacing*{\chapter}{0pt}{0pt}{1.5em}
\usepackage{xurl}
\usepackage[hidelinks]{hyperref}
\providecommand{\tightlist}{\setlength{\itemsep}{0pt}\setlength{\parskip}{0pt}}

\title{\textbf{Persona-Based Multi-Agent LLM Debate as an Automated Playtester: Validating Game Difficulty Data Against Human Performance Baselines}}
\author{Muhammad Haydar Maulana\\[0.3em] \normalsize Master of Informatics\\ \normalsize\texttt{haydariday@gmail.com}}
\date{October 2026}

\begin{document}
\hypersetup{pageanchor=false}
\maketitle
\hypersetup{pageanchor=true}
\pagenumbering{roman}
\tableofcontents
\clearpage
\pagenumbering{arabic}

\chapter{Introduction}\label{introduction}

\section{Background}\label{background}

``Is the game hard enough for the intended player?'' is one of the key
questions that game designers keep asking during development
\autocite{weber2020}. A game that is too hard drives players away while
one that is too easy soon becomes stale, both of which risk losing
players and receiving a negative reception \autocite{karavidas2022}.
That is why difficulty can be an indicator of player engagement and
experience \autocite{romero2023,alvarado2021}. Designers measure this
difficulty through playtesting. However, manual playtesting is
resource-intensive since it needs many testers over many sessions and
has to be repeated after every significant design change
\autocite{weber2020,prasetya2026,zhao2025}. For indie game developers
with limited budgets, this cost is often prohibitive. Thus, they need a
cheaper way to obtain difficulty data that still reflects human players.

Automatic playtesting has gained popularity among industry and academic
researchers alike. The methods explored so far fall into two groups.
Hand-crafted methods, such as programmed test agents
\autocite{prasetya2026} and mathematical difficulty models
\autocite{francillette2025}, require heavy engineering effort for each
game or genre. Learning-based methods such as deep reinforcement
learning need large training resources while also tending to play in
ways that humans would not \autocite{delfosse2025}. Large Language
Models (LLMs) provide an alternative, as LLM agents can play games
ranging from chess \autocite{feng2023} to Pokémon battles
\autocite{hu2024} and Minecraft \autocite{wang2023} using only a text
description of the rules. Xiao and Yang \autocite{xiao2025} show that
the difficulty LLM agents experience in Wordle and Slay the Spire
follows that of human players, even though the agents play worse.
However, these LLM playtesters rely on a single agent that plays with
only one playstyle. Human players are more diverse, as some take risks
while others play safely. The difficulty that matters to a designer is
the one experienced by this diverse population. Moreover, single LLMs
replicate human strategic behavior less closely than multi-agent systems
do \autocite{sreedhar2025}. Thus, it remains unclear whether an LLM
playtester can produce difficulty data that is both consistent and
aligned with human players once multiple playstyles are considered.

To accommodate this gap, we establish a persona-based multi-agent LLM
debate framework that acts as an automated playtester. Such a playtester
plays a game through the same actions available to a human player while
recording the outcome of each playthrough. Game difficulty here refers
to objective difficulty, which is the part of difficulty measured from
play outcomes rather than the difficulty perceived by individual players
\autocite{constant2017}. In this framework, several LLM agents with
different playstyle personas debate each in-game decision until they
reach consensus on the action to play. In other words, the agents reach
consensus on what to play rather than on how hard the game is. The
difficulty data comes from the outcome of each playthrough, such as the
number of guesses used. These outcomes are then compared with human
performance baselines, which are the average results of real players on
the same puzzles. The exact agent roles and debate protocol are defined
later in the research method.

We employ Wordle as the case study for three reasons. First, its games
are short with deterministic feedback. Second, it was studied by Xiao
and Yang \autocite{xiao2025}, so our results can be related to prior
single-agent work. Third, Wordle has public per-puzzle human
results\footnote{Engaging Data, ``Wordle Guess Distribution,''
  \url{https://engaging-data.com/wordle-guess-distribution/}.} that
express the difficulty of each puzzle as the average number of guesses
human players needed. The playtester plays a sample of historical
puzzles several times each, after which we measure how much its outcomes
vary across repeated runs of the same puzzle. We also measure how well
its average guess counts agree with the human averages, both as
continuous scores and as ordered difficulty categories. In addition, we
describe how the playstyle personas behave during the debate, since a
persona-based playtester is only meaningful when its personas actually
shape the guesses.

\section{Problem Statement}\label{problem-statement}

The background shows that game designers still lack an automated
playtester whose outcomes can serve as evidence of game difficulty for
human players. Existing LLM playtesters use a single agent with one
playstyle \autocite{xiao2025}, while the difficulty that matters to a
designer is the one experienced by a diverse player population. For a
persona-based playtester to fill this gap, two conditions must hold
before its outcomes can be used as difficulty data. First, the outcomes
must be consistent across runs. Since LLM outputs vary between calls,
the same puzzle can end differently on different runs, which leaves a
designer with little to rely on. Xiao and Yang \autocite{xiao2025} play
each Wordle puzzle 20 times yet report only the average outcome, so this
consistency is still unknown. Second, the outcomes must agree with human
difficulty \autocite{xiao2025}, because a consistent playtester that
finds every puzzle equally easy is still of no use. Moreover, the
playtester is only persona-based in practice if its personas keep their
playstyles during the debate. Debating LLM agents may conform to their
peers instead of keeping their assigned personas \autocite{baltaji2024},
which would leave a single playstyle behind the outcomes. This research
therefore examines whether the proposed playtester meets both conditions
in Wordle while describing how its personas behave.

\section{Research Questions}\label{research-questions}

Based on the problem statement, this research addresses the following
research questions:

\begin{itemize}
\tightlist
\item
  \textbf{RQ1.} How consistent are the Wordle playtest outcomes produced
  by the persona-based multi-agent debate playtester across repeated
  runs on the same puzzles?
\item
  \textbf{RQ2.} To what extent do the Wordle difficulty estimates
  derived from the playtester's outcomes agree with human difficulty,
  measured by average guess scores and difficulty categories?
\item
  \textbf{RQ3.} How do the playstyle personas behave during the debate,
  and how strongly do their suggestions shape the final guesses?
\end{itemize}

Answering RQ1 and RQ2 shows whether the playtester can be trusted as a
source of difficulty data for Wordle, which is the first step toward
solving the general problem. RQ3 shows whether the personas actually
take part in producing that data.

\section{Research Objectives}\label{research-objectives}

The objectives of this research are fourfold:

\begin{enumerate}
\def\labelenumi{\arabic{enumi}.}
\tightlist
\item
  To implement a persona-based multi-agent LLM debate playtester for
  Wordle, in which agents with different playstyle personas reach
  consensus on each guess while every playthrough is logged as
  difficulty data.
\item
  To measure the consistency of the playtester's outcomes across
  repeated runs of the same puzzles, which answers RQ1.
\item
  To measure the agreement between the playtester's difficulty estimates
  and human performance baselines using both average guess scores and
  difficulty categories, which answers RQ2.
\item
  To describe how each playstyle persona behaves during the debate and
  how often its suggestions are adopted in the final guesses, which
  answers RQ3.
\end{enumerate}

\section{Scope and Limitations}\label{scope-and-limitations}

This research focuses on Wordle as a single case study. The playtester
plays past Wordle puzzles for which public human results are available,
so that every outcome has a human counterpart to compare with. Game
difficulty is limited to objective difficulty, which is measured from
how many guesses the playtester needs to solve each puzzle
\autocite{constant2017}. The playtester is evaluated as a whole system,
judged by how consistent its outcomes are and how closely they agree
with human players. The behavior of the personas is described from the
debate records rather than tested by removing or swapping personas.

This research also has several limitations. It does not compare the
debate playtester with a single-agent playtester
\autocite{sreedhar2025}. Its persona analysis is also descriptive, so it
shows how the personas take part in the debate without proving that any
persona causes the outcomes. Moreover, the human baselines describe
players as a group rather than as individuals, so the comparison can
only be made for each puzzle as a whole. The difficulty perceived by
players is not studied either \autocite{constant2017}, since no human
participants are involved. Finally, the findings apply to Wordle and to
the LLM used in this research, while other games such as Slay the Spire
\autocite{xiao2025} may behave differently.

\section{Expected Contributions}\label{expected-contributions}

This research is expected to contribute in four ways. First, it offers
an automated playtester in which LLM agents with different playstyles
decide together how to play a game. Second, it shows whether such a
playtester gives similar results when it plays the same game again.
Third, it shows how closely the difficulty found by the playtester
follows the difficulty experienced by human players. Fourth, it shows
how the playstyle personas take part in the debate, which reveals
whether they actually shape the play.

For indie game developers, these results indicate whether an LLM
playtester can provide difficulty data that reflects human players,
which would reduce their reliance on manual playtesting
\autocite{prasetya2026,zhao2025}. For researchers, this research extends
earlier work on LLM playtesters \autocite{xiao2025} from a single agent
to several agents with different playstyles \autocite{sreedhar2025}.

\chapter{Literature Review}\label{literature-review}

\section{Automated Game Playtesting}\label{automated-game-playtesting}

Playtesting is conducted to observe how a game under development is
experienced by its players. Through playtesting, designers can find
faults in the game as well as design problems, such as levels that are
harder than intended \autocite{weber2020}. However, manual playtesting
is a resource-intensive and rigorous process, since it needs many human
testers over many sessions after every design change
\autocite{prasetya2026,zhao2025}. Automatic playtesting addresses this
cost by letting software agents play the game while recording each
playthrough.

Various methods have been explored for automatic playtesting, ranging
from hand-crafted agents to deep learning. Prasetya et al.
\autocite{prasetya2026} demonstrate a more ``classical'' yet highly
powerful approach with the APLib framework. Through goal structures and
domain-specific languages, they can create highly controlled test
agents, although this requires significant programming effort with a
steep learning curve. The mathematical model of Francillette et al.
\autocite{francillette2025} for platformer games is also noteworthy.
This model is highly deterministic and requires no sophisticated AI, yet
it is limited to the platformer genre while demanding a complex setup.

Learning-based approaches reduce part of this manual work. Gudmundsson
et al. \autocite{gudmundsson2018} train a deep neural network on human
moves in a puzzle game, after which the network predicts how hard new
levels are for players. However, such an agent can only be used once a
large set of human play data is available. Deep reinforcement learning
agents learn from their own play instead, yet they demand substantial
training resources. While these agents may outperform human players,
they lack the natural behavioral patterns of actual human gameplay
\autocite{delfosse2025}. Holmgård et al. \autocite{holmgard2019} narrow
this gap with procedural personas, which are agents tuned to represent
different playstyles. Their personas show that one playtester can cover
several ways of playing, although each persona still relies on
heuristics evolved for the specific game.

The emergence of LLMs provides an alternative approach for playtesting.
LLM agents receive the game rules and the current game state as text, so
they can play a new game without game-specific training
\autocite{xiao2025}. Zhao and Tang \autocite{zhao2025} adopt this
approach for match-3 games. They successfully address the inability of
an LLM to ``see'' the game by converting the game board into numerical
representations. As a result, their playtester achieves strong code
coverage with effective bug detection. Bateni and Whitehead
\autocite{bateni2024} show that LLM agents in Slay the Spire exhibit
long-term planning without specialized training, which lets them adapt
to changes in the game rules. Within difficulty measurement, Xiao and
Yang \autocite{xiao2025} demonstrate that LLMs can indeed be employed to
measure the difficulty of Wordle and Slay the Spire. Although their
gameplay is inferior to that of human players, the difficulty patterns
experienced by the LLMs resemble those perceived by humans.
Nevertheless, a single LLM loses too often in most games, which forces
the authors to add a handicap system. Thus, LLMs offer a promising
option for indie game developers, who rarely have the data or budget
that learning-based approaches require.

\section{Game Difficulty and Behavioral Performance
Measures}\label{game-difficulty-and-behavioral-performance-measures}

A ``fun'' game can be linked to the design of its difficulty. Players
may feel a sense of accomplishment when overcoming challenges,
particularly when those challenges are significant
\autocite{karavidas2022}. For example, Dark Souls was dubbed the hardest
game of its time, yet the higher the risk the player must take, the
greater the reward \autocite{guzsvinecz2024}. The game garnered
significant player engagement, which later gave rise to a new sub-genre
known as ``souls-like'' \autocite{klepek2017}. Given the importance of
difficulty, measuring it for the intended player is a key part of game
design.

According to Constant et al. \autocite{constant2017}, there are two
types of difficulty. First, subjective difficulty is how challenging
players perceive the game to be, which cannot be measured explicitly
because it is tied to personal skill, stress and emotion. Second,
objective difficulty is the quantifiable aspect of a game that
determines how challenging it is, independent of player perception. This
research focuses on objective difficulty only, since an automated
playtester produces play outcomes rather than perceptions of the game.

Objective difficulty is measured through behavioral performance
measures, which record what players do and achieve while playing. Such
measures include whether a challenge is completed, how many attempts it
takes and how long it takes. Dynamic difficulty adjustment systems rely
on these measures to estimate the skill of a player during play
\autocite{weber2020,romero2023}. Playtesting studies use them to compare
agents with humans, for instance the level success rate of an agent
against the pass rate of human players \autocite{gudmundsson2018}.

In Wordle, the difficulty is operationalized through the guess count,
which is how many attempts a player requires to identify the target
word. Lower guess counts indicate easier words, while higher counts or a
lost game indicate harder ones. Engaging Data\footnote{Engaging Data,
  ``Wordle Guess Distribution,''
  \url{https://engaging-data.com/wordle-guess-distribution/}.} provides
the guess distribution of real players for each past puzzle, from which
the average number of guesses per puzzle can be computed. For example, a
word with an average of 3.2 guesses is objectively easier than one with
5.8 guesses. However, the dataset does not provide additional details
such as demographic information about the players. This average serves
as the human performance baseline in this research, as it does in the
Wordle comparison of Xiao and Yang \autocite{xiao2025}.

Comparing a playtester with human players raises two separate questions.
The first is whether the playtester gives similar outcomes on the same
puzzle again, since a single run of an LLM agent may not represent its
typical play. This reliability can be assessed through outcome variance,
such as the standard deviation of repeated playthroughs. The second
question is whether the playtester ranks and groups puzzles in the same
way as human results do. Agreement can be checked on continuous scores,
such as the correlation between average guess counts, or on ordered
categories from easy to hard. Since the human baselines describe the
player population rather than individual players, both comparisons are
made per puzzle over a random sample of puzzles.

\section{Related Work}\label{related-work}

Recent studies demonstrate that LLM agents can play diverse games.
ChessGPT bridges policy learning and language modeling for chess
\autocite{feng2023}, while PokeLLMon reaches human-parity performance in
Pokémon battles \autocite{hu2024}. Voyager acts as an embodied agent in
Minecraft that keeps learning new skills as it explores
\autocite{wang2023}. However, these agents aim to play well rather than
to play like humans, which differs from the goal of playtesting.

For playtesting, the main concern is whether the results of an agent
reflect human players. The LLM playtesters reviewed above play alone
with a single playstyle \autocite{zhao2025,xiao2025}, so they do not
account for the diversity of player playstyles. Sreedhar and Chilton
\autocite{sreedhar2025} show in a strategic game that multi-agent
systems replicate human behavior more closely than single LLMs do.
Before LLMs, Holmgård et al. \autocite{holmgard2019} already represented
different players through procedural personas. More recently, MIMIC
\autocite{chen2025} equips an LLM game tester with different personality
traits, which leads to more diverse solutions with higher test coverage.
Thus, capturing the diversity of players requires moving beyond a single
LLM agent.

Multi-agent debate is one way to combine several LLM agents into a
single decision. In a debate, agents propose their own answers before
reading the reasoning of the others, after which they revise their
answers over several rounds until they agree \autocite{du2024}. Du et
al. \autocite{du2024} show that this process improves factual accuracy
and reasoning compared with a single agent. Liang et al.
\autocite{liang2024} let agents express arguments in a state of ``tit
for tat'', where each agent directly responds to the position of the
previous agent. Meanwhile, a managing agent oversees the exchange to
obtain the final decision. ChatEval \autocite{chan2024} assigns
different roles to agents in a debate for evaluating text, while finding
that the diversity of these roles is important for its results. However,
Smit et al. \autocite{smit2024} find that debate does not reliably
outperform other prompting strategies, as its results depend strongly on
how the debate is set up. Personas also do not always hold in a debate,
since agents may conform to their peers or drift away from their
assigned views \autocite{baltaji2024}.

We adopt the debate structure of Liang et al. \autocite{liang2024} for
two reasons. First, its managing agent provides a single point where the
arguments end in one action, which a playtester needs because every turn
must produce a playable move. Second, agents that argue from different
positions keep the debate from reaching agreement immediately, which
matches the role of playstyle personas \autocite{chan2024}. Since the
benefit of debate depends on its setup \autocite{smit2024}, we evaluate
the resulting playtester directly rather than assuming that debate
improves it.

\section{Research Gap and
Positioning}\label{research-gap-and-positioning}

The reviewed studies reveal a trade-off between generality and human
likeness. Hand-crafted playtesters are highly controlled but must be
rebuilt for each game \autocite{prasetya2026,francillette2025}.
Learning-based playtesters require either large sets of human play data
or substantial training resources
\autocite{gudmundsson2018,delfosse2025}. LLM playtesters can play a new
game from its rules alone, while the difficulty they experience can
follow that of human players \autocite{xiao2025}. However, existing LLM
playtesters rely on a single agent with one playstyle, even though
personas have long been used to represent different players
\autocite{holmgard2019}.

While multi-agent debate could reconcile these heterogeneous playstyles
into one consensus action, automated consensus frameworks remain absent
from current difficulty measurement methods. Thus, three questions
remain open. First, LLM outputs vary between calls, so it is unclear
whether a playtester built from several debating agents produces
consistent outcomes across repeated runs. Even Xiao and Yang
\autocite{xiao2025} play each Wordle puzzle 20 times yet report only the
average outcome. Second, it is unclear whether the difficulty found by
such a playtester still agrees with human performance baselines. Third,
it is unclear whether the personas keep their playstyles during the
debate, since debating agents may conform to their peers
\autocite{baltaji2024}. To the best of our knowledge, no prior study has
applied persona-based multi-agent debate to measure game difficulty.

To accommodate this gap, our research is positioned where these lines of
work meet. We take playstyle personas from persona-based playtesting
\autocite{holmgard2019} while taking the consensus process from
multi-agent debate \autocite{liang2024}. Both are then applied to the
LLM difficulty measurement of Xiao and Yang \autocite{xiao2025}, using
Wordle and its public human results. Unlike studies that compare
multi-agent systems with a single agent
\autocite{sreedhar2025,smit2024}, this research validates the playtester
as a system. It does so through the consistency of its outcomes (RQ1),
their agreement with human difficulty (RQ2) and the behavior of its
personas during the debate (RQ3).

\printbibliography[heading=bibintoc,title={References}]
\end{document}
